\section{Mixed-wake interleaving at arbitrary width}
\label{app:mixed}
The thin-wake classification uses a particular growth and width. A separate
parity identity constrains arbitrary finite odd-class wakes at the same full
drift \((2,-2)\), without a period or width cutoff. It tests the transverse
order of the classes, in addition to their signed residue counts.

Translate a vertical-arrival phase to the origin. Then even levels
\(t=x+y\) have vertical arrivals, denoted \(Q\), and odd levels have
horizontal arrivals, denoted \(C\). Lanes are \(r=x\bmod2\). Let \(O\)
be the odd classes, \(g=|O|\), and define
\[
 V=|O\cap Q|,\qquad
 J(O)=\#\{(q,c)\in(O\cap Q)\times(O\cap C):t(q)<t(c)\}.
\]
Distinct classes on the same layer are counted separately.

\begin{theorem}[Mixed-wake interleaving]\label{thm:mixed}
Every finite-support periodic realization of full drift \((2,-2)\) satisfies
\[
 V\equiv0\pmod2,\qquad J(O)+V/2\equiv1\pmod2.
\]
The condition can be evaluated in \(O(g\log g)\) integer operations by sorting
only the occupied odd classes, irrespective of transverse gap sizes.
\end{theorem}

\subsection{Signed residues and the rank factorization}
We use the stable-class and signed-residue arguments of
Jillhewar~\cite{Jillhewar2026}. The supplemental
\texttt{proof/periodic-growth-closure.md} supplies an attributed derivation,
including an explicit mod-four potential. For a function
\(f:\Z\to\Z/4\Z\), choose \(S(0)=0\) and
\(S(x+1)-S(x)=f(x)+f(x+1)\). With \(\chi=(-1)^{x+y}\), put
\[
\begin{aligned}
 F_N&=-\chi S(x),&F_E&=\chi(S(x)-f(x)),\\
 F_S&=-\chi(S(x)+2f(x)),&F_W&=\chi(S(x)+f(x)).
\end{aligned}
\]
The complete-state quantity
\(\sum_{(x,y)\in B}\chi(x,y)f(x)+F_h(p)\) is conserved modulo four,
as direct substitution in the eight heading/colour cases verifies.
For a residue indicator modulo a nonzero horizontal displacement \(a\),
its monodromy gives
\[
 \sum_{q\in O,\ q_x=r\bmod|a|}\chi(q)\equiv2\pmod4.
\]
Quarter-turning gives the vertical counterpart. In particular, with
\(q_r=|O\cap Q\text{ in lane }r|\) and
\(c_r=|O\cap C\text{ in lane }r|\),
\[
 q_0-c_0\equiv q_1-c_1\equiv q_0-c_1\equiv q_1-c_0\equiv2\pmod4.
\]
Thus \(V\) is even and \(V/2\equiv q_0\pmod2\).

For completeness we state the port argument used in the new identity.
At each class, order visits by stabilized rank \(k\), and record positive
incoming and outgoing port bits \(\alpha_k,\beta_k\): the positive ports
are N/E at \(Q\), E/N at \(C\). For a binary string with \(m\) positions
and \(A\) ones,
\[
 \operatorname{inv}(\alpha)
 \equiv A(m-1)+\sum_k k\alpha_k+\binom A2\pmod2.
\]
Apply this to both port strings. Directed-edge binomial terms cancel between
source and target. If \(D_v\) is the odd-class bit and
\(p_v=\sum_k(\alpha_k+\beta_k)\bmod2\), the remaining local correction is
\(p_v(D_v+1)\) at \(Q\), and \(D_v(p_v+1)\) at \(C\).
Denote its sum by \(\mathcal C\), and the edge-block permutation sign by
\(\mathcal E\). Stable order within each directed edge block is preserved,
since its successor time changes by the fixed amount \(1-P\,\mathrm{voltage}\).
The phase successor is one even-length cycle. Its factorization consequently
gives
\[
 K:=\sum_{\text{L events}}\mathrm{rank}
 \equiv1+\mathcal E+\mathcal C\pmod2.
\]
For an actual stable class, L has odd rank. Also \(4\mid P\) and
\(4\mid g\), so \(K\equiv(P-g)/2\equiv0\pmod2\). Hence
\(\mathcal E+\mathcal C=1\). This factorization is inherited; the following
evaluation for general mixed wakes is the extension.

\subsection{A Boolean identity with a telescoping boundary term}
All arithmetic in this subsection is over \(\mathbb F_2\). Write
\(Q_{j,r}=D_{(2j,r)}\), \(C_{j,r}=D_{(2j+1,r)}\). Let \(e,w\) count
E/W outgoing edges at \(Q\) modulo two, and \(n,s\) count N/S outgoing
edges at \(C\). Local balance gives, as two-vector identities,
\[
 w_j=e_j+Q_j,\quad s_j=n_j+C_j,\quad
 n_j=n_{j-1}+Q_j+C_j,\quad
 e_j=e_{j-1}+Q_j+\operatorname{swap}(C_{j-1}).
\]
All counts vanish below and above the finite range of one period.
We prove the closed identity
\[
 \mathcal E+\mathcal C=J(O)+\binom g2+q_0.                 \tag{*}
\]

Encode a layer's outgoing parities in the order negative lane zero,
positive lane zero, negative lane one, positive lane one. Its internal
inversion contributions are
\[
 S_Q(b)=b_0b_2+b_1b_2+b_1b_3,\qquad S_C(b)=b_1b_2.
\]
The adjacent-layer term is \(A(a,b)=(a_1+a_3)(b_0+b_2)\). At separation
two the terms are
\[
 B_Q(a,b)=a_1(b_0+b_2)+a_3b_2,\qquad
 B_C(a,b)=a_1b_0+a_3(b_0+b_2).
\]
Larger separations give no inversions. These expressions follow by ordering
the directed targets: Q edges swap lane and C edges preserve it.

Before a new pair of layers, use state \((e_0,e_1,c_0,c_1,z)\), where
\(z\) is the accumulated Q population parity. The previous N parities are
\(n_0=e_1+z+c_0\), \(n_1=e_0+z+c_1\). For new bits
\((Q_0,Q_1,C_0,C_1)\), the next state is
\((e'_0,e'_1,C_0,C_1,z')\), where
\[
 e'_0=e_0+Q_0+c_1,\quad e'_1=e_1+Q_1+c_0,\quad
 n'_r=n_r+Q_r+C_r,\quad z'=z+Q_0+Q_1.
\]
Set
\[
\begin{aligned}
 a&=(0,e_0,0,e_1),& b&=(n_0+c_0,n_0,n_1+c_1,n_1),\\
 u&=(e'_0+Q_0,e'_0,e'_1+Q_1,e'_1),&
 v&=(n'_0+C_0,n'_0,n'_1+C_1,n'_1).
\end{aligned}
\]
Then
\[
\begin{aligned}
 \Delta\mathcal E&=S_Q(u)+S_C(v)+A(b,u)+B_Q(a,u)+A(u,v)+B_C(b,v),\\
 \Delta\mathcal C&=(n_0+e'_0)(1+Q_0)+(n_1+e'_1)(1+Q_1)\\
 &\quad+C_0(1+e'_1+n'_0)+C_1(1+e'_0+n'_1),\\
 \Delta J&=z'(C_0+C_1).
\end{aligned}
\]
For \(d=(Q_0,Q_1,C_0,C_1)\), the population correction is
\[
 \Delta\binom g2=(n_0+n_1)\sum_i d_i+\sum_{i<j}d_id_j.
\]
The explicit boundary polynomial
\[
 H=e_0+e_1+z+e_0c_1+e_1c_0+e_1c_1+e_1z+c_0c_1+c_0z
\]
satisfies
\[
 H(\mathrm{old})+H(\mathrm{new})=
 \Delta\mathcal E+\Delta\mathcal C+\Delta J+
 \Delta\binom g2+Q_0.
\]
Substitution with \(x^2=x\) proves the identity in nine Boolean variables;
the checker both expands the polynomial and tests all 512 assignments.
Summing over any finite number of layer pairs telescopes \(H\). Empty
boundary padding makes it zero at both ends, giving (*). Since \(4\mid g\)
and \(q_0\equiv V/2\), comparison with
\(\mathcal E+\mathcal C=1\) proves Theorem~\ref{thm:mixed}.

\subsection{Growth-eight wakes with two full extreme layers}
Suppose \(g=8\) and both residue classes at each transverse extreme are odd.
Extreme visits are singleton right turns, hence on C layers. They supply four
C classes. Thus \(V\) is even and at most four. If \(V=0\), the new parity
is zero. If \(V=4\), all C classes are at the extremes; each Q lies below
exactly two of them, so \(J=8\) and \(J+V/2\) is again even. Therefore
\(V=2\), with lane populations \((q_0,q_1)=(1,1)\),
\((c_0,c_1)=(3,3)\). The x/y residue matrix is
\(\left[\begin{smallmatrix}1&3\\3&1\end{smallmatrix}\right]\).

If the four internal odd classes have distinct levels, their increasing
Q/C sequence must have an even number of Q-before-C pairs. The alternating
orders QCQC and CQCQ are excluded for arbitrary gaps and lane choices.
QQCC, QCCQ, CQQC and CCQQ pass this particular test; they are candidates
requiring the full realization criterion. This consequence does not require
a transverse-width bound.

The exact controls include 18,895 geometric heading-restoring words,
65,536 finite layer subsets, all 104 standard phases, and an independent
edge-permutation computation. A separate bounded screen at widths 10,12,14
excludes respectively 3,058, 11,374 and 33,320 assignments passing the signed
residues and checker-mixing filters. Those finite tables illustrate the
screen; the telescoping identity proves it at every width. The complete
periodic-growth supplement also retains attributed small-drift exclusions
and the regenerated unit-diagonal certificate, keeping their ownership and
logical dependencies explicit.
